% GENERATED FILE. Edit canonical lessons, then run npm run generate:books.
% canonical-source-hash: fnv1a64:6c788bc0104707ff
% canonical-lessons: TA-C29-kaalai-vanakkam

\chapter{Good Morning}
\label{ch:ta-good-morning}

This chapter is generated from the canonical micro-lessons used by Language Ladder.
Each section stays independently resumable and preserves the authored prerequisite order.

\section[\ta{காலை} \ta{வணக்கம்}]{\ta{காலை} \ta{வணக்கம்} — "morning greetings"}

\label{lesson:TA-C29-kaalai-vanakkam}



\begin{quote}
\textbf{Warm-up.} {[}PAUSE 2s{]} Every other language in this arc built its "good morning" on a word meaning "good" or "auspicious." Tamil does something completely different: it just puts its everyday greeting word after its word for "morning."
\end{quote}

\subsection*{You'll want to know: \ta{காலை} \ta{வணக்கம்} — literally, "morning greetings"}
\textbf{\ta{காலை} \ta{வணக்கம்}} (\textbf{kālai vaṇakkam}) — "\textbf{good morning}" — is confirmed, by two independently-fetched sources, as Tamil's standard greeting for this time of day, used "\textbf{from sunrise until around 11 AM}." It's built from two words you already have: \textbf{\ta{காலை}} (already met, Chapter 26, "\textbf{morning}") plus \textbf{\ta{வணக்கம்}} (already met, Chapter 1, "\textbf{reverence, homage}" — Tamil's everyday word for "hello"). Literally: "\textbf{morning salutations}."

\begin{culture}
Here's the honest structural surprise: Hindi, Kannada, Telugu, and Malayalam each build their morning greeting on a \textbf{śubha}-equivalent word ("good, auspicious") plus a time-word. Tamil does something genuinely \textbf{different} — there's no "good/auspicious" word here at all. It's simply \textbf{\ta{வணக்கம்}}, Tamil's own general greeting, specified with \textbf{\ta{காலை}} in front of it to mean "morning" greetings specifically. Also be honest about a claim that didn't hold up: an initial search summary suggested a Sanskrit-borrowed word, "suprabhatham," serves as a casual alternative — but that claim appears on \textbf{neither} of the two pages actually fetched, so it's deliberately \textbf{not} asserted here.
\end{culture}

\subsection*{Guided Practice}
{[}PAUSE 1s{]}

\begin{itemize}
\raggedright
  \item {[}YOU SAY: "kālai vaṇakkam" — "good morning," literally "morning greetings"{]}
  \item {[}YOU SAY: the honest structural surprise — no śubha-equivalent word here at all, unlike every other language in this arc{]}
  \item {[}YOU SAY: the two pieces you already know — kālai (Chapter 26) + vaṇakkam (Chapter 1){]}
\end{itemize}

\begin{tcolorbox}[breakable,colback=teal!4,colframe=teal!35!black,title={Wrap-up recall}]
{[}PAUSE 3s{]}

\begin{itemize}
\raggedright
  \item What two known words build \ta{காலை} \ta{வணக்கம்}? (\textbf{\ta{காலை}}, "morning" + \textbf{\ta{வணக்கம்}}, "reverence/hello.")
  \item Does it follow the śubha-equivalent + time-word pattern used by Hindi, Kannada, Telugu, and Malayalam? (\textbf{No} — \ta{காலை} simply specifies Tamil's own general greeting, \ta{வணக்கம்}.)
  \item Is "suprabhatham" confirmed as a casual alternative? (\textbf{No} — neither fetched source supports that claim, so this lesson does not assert it.)
\end{itemize}
\end{tcolorbox}
